% ==================== 补充材料 ====================
\documentclass[12pt,a4paper]{article}
\usepackage[utf8]{inputenc}
\usepackage{graphicx}
\usepackage{booktabs}
\usepackage{geometry}
\usepackage{xeCJK}
\setCJKmainfont{SimSun}
\geometry{left=2.54cm,right=2.54cm,top=2.54cm,bottom=2.54cm}
\onehalfspacing

\title{\heiti\Large 补充材料\\[0.5em]\normalsize 肺癌靶向治疗耐药机制研究}
\date{\today}

\begin{document}
\maketitle
\tableofcontents
\newpage

\section{补充表 1：EGFR 突变类型与临床特征}

\begin{table}[h]
\centering
\begin{tabular}{p{3cm}p{2cm}p{2cm}p{4cm}}
\toprule
\textbf{突变类型} & \textbf{外显子} & \textbf{频率} & \textbf{临床特征} \\
\midrule
外显子 19 缺失 & 19 & 45\% & 亚洲人、女性、非吸烟者多见 \\
L858R 点突变 & 21 & 40\% & 腺癌、女性多见 \\
G719X & 18 & 2-3\% & 二代 TKI 更优 \\
S768I & 20 & 1-2\% & 常与其他突变共存 \\
L861Q & 21 & 1-2\% & 二代 TKI 更优 \\
外显子 20 插入 & 20 & 2-3\% & 原发耐药 \\
T790M & 20 & 60\% & 获得性耐药 \\
C797S & 20 & 15-20\% & 三代 TKI 耐药 \\
\bottomrule
\end{tabular}
\end{table}

\newpage

\section{补充表 2：EGFR-TKI 临床试验汇总}

\begin{table}[h]
\centering
\begin{tabular}{p{2cm}p{2cm}p{2.5cm}p{2cm}p{2.5cm}}
\toprule
\textbf{试验} & \textbf{药物} & \textbf{人群} & \textbf{中位 PFS} & \textbf{参考文献} \\
\midrule
IPASS & 吉非替尼 & EGFR+ 一线 & 9.5 vs 6.3 月 & NEJM 2009 \\
EURTAC & 厄洛替尼 & EGFR+ 一线 & 9.7 vs 5.2 月 & Lancet Oncol 2012 \\
LUX-Lung 3 & 阿法替尼 & EGFR+ 一线 & 11.1 vs 6.9 月 & JCO 2013 \\
FLAURA & 奥希替尼 & EGFR+ 一线 & 18.9 vs 10.2 月 & NEJM 2018 \\
FLAURA2 & 奥希替尼+化疗 & EGFR+ 一线 & 25.5 vs 16.7 月 & NEJM 2023 \\
\bottomrule
\end{tabular}
\end{table}

\newpage

\section{补充表 3：ALK 融合变异类型}

\begin{table}[h]
\centering
\begin{tabular}{p{3cm}p{3cm}p{2cm}p{4cm}}
\toprule
\textbf{变异类型} & \textbf{结构} & \textbf{频率} & \textbf{临床特征} \\
\midrule
V1 (E13;A20) & EML4 外显子 1-13 + ALK 外显子 20-29 & 33\% & 最常见 \\
V2 (E20;A20) & EML4 外显子 1-20 + ALK 外显子 20-29 & 7\% & 较少见 \\
V3 (E6a/b;A20) & EML4 外显子 1-6 + ALK 外显子 20-29 & 28\% & 可能预后较差 \\
其他融合伙伴 & KIF5B、TFG、HIP1 等 & 32\% & 非 EML4 融合 \\
\bottomrule
\end{tabular}
\end{table}

\newpage

\section{补充表 4：ALK 抑制剂临床试验}

\begin{table}[h]
\centering
\begin{tabular}{p{2cm}p{2cm}p{2.5cm}p{2.5cm}p{2cm}}
\toprule
\textbf{试验} & \textbf{药物} & \textbf{人群} & \textbf{中位 PFS} & \textbf{参考文献} \\
\midrule
PROFILE 1014 & 克唑替尼 & ALK+ 一线 & 10.9 vs 7.0 月 & NEJM 2014 \\
J-ALEX & 阿来替尼 & ALK+ 一线 & NR vs 10.2 月 & Lancet 2017 \\
ALEX & 阿来替尼 & ALK+ 一线 & 34.8 vs 10.9 月 & NEJM 2017 \\
CROWN & 洛拉替尼 & ALK+ 一线 & NR vs 9.3 月 & NEJM 2020 \\
ALTA-1L & 布格替尼 & ALK+ 一线 & 24.0 vs 9.2 月 & JCO 2020 \\
\bottomrule
\end{tabular}
\end{table}

\newpage

\section{补充表 5：耐药机制频率汇总}

\begin{table}[h]
\centering
\begin{tabular}{p{3cm}p{5cm}p{4cm}}
\toprule
\textbf{靶点} & \textbf{耐药机制} & \textbf{频率} \\
\midrule
\multirow{6}{*}{EGFR} & T790M 突变 & 50-60\% (一代/二代 TKI) \\
& C797S 突变 & 15-20\% (奥希替尼) \\
& MET 扩增 & 15-20\% \\
& HER2 扩增/突变 & 5-7\% \\
& 小细胞转化 & 5-10\% \\
& 其他/未知 & 10-15\% \\
\midrule
\multirow{4}{*}{ALK} & ALK 激酶域突变 & 30-50\% \\
& \quad - G1202R & 10-30\% \\
& 旁路激活 & 20-30\% \\
& 未知 & 20-30\% \\
\midrule
\multirow{3}{*}{ROS1} & ROS1 激酶域突变 & 40-50\% \\
& \quad - G2032R & 40\% \\
& 旁路激活 & 30-40\% \\
\bottomrule
\end{tabular}
\end{table}

\newpage

\section{补充表 6：液体活检检测方法}

\begin{table}[h]
\centering
\begin{tabular}{p{2.5cm}p{2.5cm}p{2.5cm}p{2.5cm}p{3cm}}
\toprule
\textbf{方法} & \textbf{灵敏度} & \textbf{特异性} & \textbf{检测限} & \textbf{优势} \\
\midrule
ddPCR & 高 & 高 & 0.01\% & 快速、定量 \\
NGS (Panel) & 高 & 高 & 0.1\% & 多基因检测 \\
NGS (WES) & 中 & 高 & 1\% & 全外显子 \\
ARMS-PCR & 中 & 高 & 1\% & 临床常用 \\
\bottomrule
\end{tabular}
\end{table}

\newpage

\section{补充表 7：在研新一代抑制剂}

\begin{table}[h]
\centering
\begin{tabular}{p{2cm}p{2.5cm}p{2.5cm}p{3cm}p{2.5cm}}
\toprule
\textbf{药物} & \textbf{靶点} & \textbf{研发阶段} & \textbf{适应症} & \textbf{特性} \\
\midrule
BLU-945 & EGFR (四代) & Phase I/II & C797S 突变 & CNS 渗透 \\
BBT-176 & EGFR (四代) & Phase I & C797S 突变 & 自主研发 \\
Zotizalkib & ALK (三代) & Phase II & G1202R & CNS 渗透 \\
Repotrectinib & ROS1 (二代) & Approved & G2032R & FDA 2023 批准 \\
\bottomrule
\end{tabular}
\end{table}

\newpage

\section{补充方法 1：文献检索策略}

\subsection*{检索数据库}
\begin{itemize}
    \item PubMed/MEDLINE
    \item Embase
    \item Cochrane Library
    \item Web of Science
    \item ClinicalTrials.gov
\end{itemize}

\subsection*{检索词}
\textbf{EGFR 相关}：
\begin{verbatim}
("lung neoplasms" OR NSCLC) AND ("EGFR") 
AND ("tyrosine kinase inhibitor" OR TKI)
AND ("drug resistance" OR "resistance mechanism")
\end{verbatim}

\textbf{ALK 相关}：
\begin{verbatim}
("lung neoplasms" OR NSCLC) AND ("ALK")
AND ("tyrosine kinase inhibitor" OR TKI)
AND ("drug resistance" OR "resistance mechanism")
\end{verbatim}

\subsection*{纳入标准}
\begin{enumerate}
    \item 研究对象：经病理确诊的 NSCLC 患者
    \item 研究类型：RCT、队列研究、病例系列、综述
    \item 语言：英语或中文
    \item 发表日期：2010 年 1 月 1 日至 2026 年 3 月 31 日
\end{enumerate}

\subsection*{排除标准}
\begin{enumerate}
    \item 个案报道（n < 5）
    \item 会议摘要（无全文）
    \item 动物研究
    \item 重复发表
\end{enumerate}

\newpage

\section{补充讨论 1：研究局限性}

\begin{enumerate}
    \item \textbf{异质性}：纳入研究的患者特征、检测方法存在异质性
    \item \textbf{发表偏倚}：阳性结果更容易发表
    \item \textbf{检测方法的演变}：早期 PCR，晚期 NGS，灵敏度不同
    \item \textbf{液体活检的局限性}：可能遗漏某些耐药机制
    \item \textbf{时间因素}：耐药机制随治疗时间演变
    \item \textbf{未知机制}：约 10-30\% 的耐药机制仍不明确
\end{enumerate}

\newpage

\section{补充讨论 2：临床实践建议}

\subsection*{基线检测}
\begin{itemize}
    \item 所有晚期 NSCLC 患者应进行 EGFR、ALK、ROS1 检测
    \item 优先使用 NGS panel
    \item 组织样本优先，液体活检作为补充
\end{itemize}

\subsection*{进展后检测}
\begin{itemize}
    \item 疾病进展时应进行重复分子检测
    \item 优先组织活检，无法获取时使用液体活检
    \item 检测范围应覆盖已知耐药机制
\end{itemize}

\subsection*{治疗选择}
\begin{itemize}
    \item 根据耐药机制选择针对性治疗
    \item 鼓励参与临床试验
    \item 考虑联合治疗策略
\end{itemize}

\subsection*{随访监测}
\begin{itemize}
    \item 定期影像学评估（每 2-3 个月）
    \item 动态 ctDNA 监测可早期发现耐药
    \item 关注 CNS 转移（定期脑部 MRI）
\end{itemize}

\end{document}
